\subsection{Energy-based Learning}
Energy-based learning (EBL) \citep{lecun2006tutorial,dawid2024introduction} provides a unifying framework for modeling prediction as minimizing an energy function. Early forms include Hopfield networks \citep{hopfield1982neural} and Boltzmann machines \citep{ackley1985learning}. Modern developments in EBL span both generative modeling—via energy functions \citep{du2019implicit} or their gradients (as in score-based models \citep{pmlr-v37-sohl-dickstein15,song2019generative})—and representation learning. Another line of work views network layers as the result of iterative energy minimization. Some approaches define energy implicitly through neural networks \citep{bai2019deep,du2022learning, Du_2024_ICML}, while recent work Energy Transformer \citep{hoover2024energy} draws analogies between attention layers and explicit energy descent but mainly focuses on reinterpretation rather than principled derivation. Our work differs in that we design the Transformer block by quantifying a general principle that can induce variants through alternative energy.

Other works also explore energy formulations on the hypersphere \citep{liu2018learning,loshchilov2025ngpt}, but mostly in the weight space. By contrast, we define our energies directly on the representation space. Additionally, recent theoretical studies on memory capacity in modern Hopfield networks \citep{hu2024computational,pmlr-v235-wu24i,hu2024provably} emphasize spreading patterns on the sphere but focus primarily on memory retrieval and cross-attention. 

\subsection{Model Design from First Principles}
While neural network architectures are often shaped by engineering practices, recent work has explored designing or interpreting them through principled lenses like signal processing, information theory, and neurobiology. For example, deep unrolling of the sparse coding algorithms has led to the development of fully connected networks \citep{gregor2010learning}, convolution networks \citep{papyan2017convolutional, papyan2018theoretical}, and even graph neural networks through iterative algorithms \citep{yang2021graph}. Similarly, the sparse rate reduction principle has been used to derive the Transformer architecture \citep{yu2023white}. Other approaches draw inspiration from approximation theory \citep{liu2025kan} and brain computation \citep{kozachkov2023building}, further bridging the gap between theoretical insights and practical network design. 

\subsection{Depth Recurrence in Transformers}
Recurrence has been scientifically considered as a core computational mechanism in the biological visual system enabling flexible computational depth, integration of priors, and efficient allocation of compute under resource constraints \citep{van2020going}. Depth-wise recurrence in Transformer architectures, under the name of universal \citep{devlin-etal-2019-bert}, recursive/recurrent \citep{bae2025relaxed,geiping2025scaling} or looped Transformers \citep{giannou2023looped}, where cross-layer weights are shared, has emerged as a critical avenue for reducing parameters and enabling iterative reasoning. Its repeated reuse of the same layer has been demonstrated to be capable of emulating iterative algorithms \citep{schwarzschild2021can,saunshi2025reasoning}. Crucially, it has been demonstrated with particular strength mostly on systematic generalization \citep{csordas2021devil} and structured reasoning tasks \citep{schwarzschild2021can,bansal2022end}, while its applications on general-domain tasks that need hierarchical processing (\eg, image perception) are under-explored. 
